% Vorlage fuer Abschlussarbeiten am Fachbereich Mathematik der TU Darmstadt.
% Geeignet fuer Bachelorarbeiten, Masterarbeiten und Diplomarbeiten.

\documentclass[12pt,a4paper,twoside]{report}
% Benutzt man die Option draft, so kann man die Umbrueche ueberpruefen.

% Hier werden alle benoetigten Pakete und Einstellungen geladen. Auch
% hier sind Sie frei diese direkt in der Praeambel zu laden.


\usepackage[ngerman]{babel}

% Kodierung fuer unixoide Systeme und Windowssysteme.

\usepackage[utf8]{inputenc}

% Schriftart Times New Roman.

\usepackage{times}

% Erweitert den Zeichenvorrat, so dass z.B. auch Umlaute im PDF-Dokument
% gefunden werden.

\usepackage[T1]{fontenc}

% Zur Einbindung von Bildern.

\usepackage{graphicx}

% Erweiterte enumerate-Umgebung.

\usepackage{enumerate}

% Verschiedene Pakete, die nuetzlich sind um Mathematik in LaTeX zu setzen.

\usepackage{amsmath, amssymb, amsthm, dsfont}

% Indentation

\usepackage{scrextend}

% Spaces for newcommand

\usepackage{xspace}

% tikz
\usepackage{tikz}
%\usetikzlibrary{arrows.meta}% arrow tip library
%\usetikzlibrary{bending}% better arrow head for bended lines

% Hyphenation of wordas that already contain a hyphen
% http://tex.stackexchange.com/questions/2706/adequate-hyphenation-of-words-already-containing-a-hyphen
% Example:
% \textsc{Alexandroff}\hyp{}\textsc{Urysohn}\hyp{}Kompaktheit

\usepackage{hyphenat}

% Hyperlinks

%\usepackage{hyperref}
\usepackage[pdfauthor={Fabian Gabel},%
pdftitle={Bachelorthesis - Helmholtz-Zerlegung},%
%pagebackref=true,%
%pdftex
]{hyperref}


\usepackage[
%nonumberlist, %keine Seitenzahlen anzeigen
toc,          %Einträge im Inhaltsverzeichnis
section]      %im Inhaltsverzeichnis auf section-Ebene erscheinen
{glossaries}

%%%%%%%%%%%%%%%%
% Seitenlayout %
%%%%%%%%%%%%%%%%

% DIV# gibt den Divisor für die Layoutberechnung an.
% Vergrößern des Divisors vergrößert den Textbereich.
% BCOR#cm gibt die Breite des Bundstegs an.
\usepackage[DIV14,BCOR2cm]{typearea}

% Abstand obere Blattkante zur Kopfzeile ist 2.54cm - 15mm
\setlength{\topmargin}{-15mm}

% Keine Einrueckung nach einem Absatz.

\parindent 0pt

% Abstand zwischen zwei Abs\"atzen.

\parskip 12pt

% Zeilenabstand.

\linespread{1.25}

% Inhaltsverzeichnis erstellen.

\usepackage{makeidx}
\makeindex


% Hier werden Makros und Umgebungen eingebunden. Diese werden separat
% in der Datei befehle.tex definiert. Sie sind frei diese Befehle auch
% direkt in der Praeambel zu definieren.
% Makros

\newcommand{\N}{\mathbb{N}} % natuerliche Zahlen
\newcommand{\Z}{\mathbb{Z}} % ganze Zahlen
\newcommand{\Q}{\mathbb{Q}} % rationale Zahlen
\newcommand{\R}{\mathbb{R}} % reelle Zahlen
\newcommand{\K}{\mathbb{K}} % Körper
\newcommand{\co}{\mathrm{co}} % compact open
\newcommand{\W}{\mathcal{W}} % Weyl-Gruppe

\DeclareMathOperator{\Ret}{Re} % Realteil
%\newcommand{\PIT}{\textbf{PIT }} % Prime Ideal Theorem
\newcommand{\AC}{\textbf{AC}\xspace} % Prime Ideal Theorem
\newcommand{\UFT}{\textbf{UFT}\xspace} % Prime Ideal Theorem
\newcommand{\PIT}{\textbf{PIT}\xspace}
\newcommand{\ZF}{\textbf{ZF}\xspace}
\newcommand{\ZFC}{\textbf{ZFC}\xspace}
\newcommand{\CC}{\textbf{CC}\xspace}
\newcommand{\CCR}{\textbf{CC}($\R$)\xspace}
\newcommand{\PCC}{\textbf{PCC}\xspace}
\newcommand{\PCCR}{\textbf{PCC}($\R$)\xspace}

\newcommand{\grad}[1]{\gv{\nabla} #1} % for gradient
\let\divsymb=\div % rename builtin command \div to \divsymb
\renewcommand{\div}[1]{\mathrm{div} #1} % for divergence
\renewcommand{\d}[1]{\ensuremath\, {\operatorname{d}\!{#1}}}

\newcommand{\banach}{\textsc{Banach}}
\newcommand{\cauchy}{\textsc{Cauchy}}
\newcommand{\euklid}{\textsc{Euklid}}
\newcommand{\fubini}{\textsc{Fubini}}
\newcommand{\hahn}{\textsc{Hahn}}
\newcommand{\helmholtz}{\textsc{Helmholtz}}
\newcommand{\hoelder}{\textsc{Hölder}}
\newcommand{\killing}{\textsc{Killing}}
\newcommand{\lebesgue}{\textsc{Lebesgue}}
\newcommand{\leibniz}{\textsc{Leibniz}}
\newcommand{\lie}{\textsc{Lie}}
\newcommand{\lipschitz}{\textsc{Lipschitz}}
\newcommand{\poincare}{\textsc{Poincaré}}
\newcommand{\sobolev}{\textsc{Sobolev}}
\newcommand{\weyl}{\textsc{Weyl}}


% Umgebungen für Definitionen, Sätze, usw.

%\newtheorem{defn}{Definition}[section]
%\newtheorem{ex}{Beispiel}[chapter]

\theoremstyle{plain}
\newtheorem{thm}{Satz}[chapter]
\newtheorem{lem}[thm]{Lemma}
\newtheorem{cor}[thm]{Korollar}
\newtheorem{prop}[thm]{Proposition}

\theoremstyle{definition}
\newtheorem{defn}[thm]{Definition}

\theoremstyle{remark}
\newtheorem*{bem}{Bemerkung}

\def\Satzrefname{Satz}

\DeclareMathOperator{\spn}{span}
\DeclareMathOperator{\sign}{sign}
\DeclareMathOperator{\GL}{GL}
\DeclareMathOperator{\Sym}{Sym}
\DeclareMathOperator{\hgt}{ht}
\DeclareMathOperator{\dist}{dist}
\DeclareMathOperator{\loc}{loc}

%\renewcommand{\refname}{Literaturverzeichnis}
\addto\captionsngerman{%
\renewcommand*{\bibname}{Literaturverzeichnis}
}

%\makeglossaries

\pagestyle{fancy}
\fancypagestyle{MyStyle}{%
\fancyhf{} %Clean headers
\fancyhead[LE,RO]{\thepage}
\fancyhead[RE]{\normalfont \nouppercase{\leftmark}}
\fancyhead[LO]{\normalfont \nouppercase{\rightmark}}
\fancyfoot[C]{\tiny Version vom \today}
\renewcommand{\chaptermark}[1]{\markboth{\thechapter\ {\normalfont{##1}}}{}}
\renewcommand{\sectionmark}[1]{\markright{\thesection\ {\normalfont{##1}}}{}}
}

\fancypagestyle{plain}{%
\fancyhf{} % clear all header and footer fields
\fancyfoot[LE,RO]{\thepage} % except the center
\renewcommand{\headrulewidth}{0pt}
\renewcommand{\footrulewidth}{0pt}}


\begin{document}
% Auf der Titelseite und im Inhaltsverzeichnis sollen keine
% Seitenzahlen dargestellt werden.
\pagestyle{empty}

% Hier wird die Titelseite eingebunden.
\begin{titlepage}
  \begin{center}
    \vspace{1cm}
    \includegraphics[width=0.5\linewidth]{TU_Darmstadt_Logo.pdf}
    \vspace{1cm}
    
    \large{Fachbereich Mathematik}
    \vspace{2.5cm}
    
    \large{Bachelorarbeit}
    \vspace{2cm}

    \Huge{Die Helmholtz-Zerlegung in $L^2$}
    
    \vspace*{3cm}    
    
		\large
                \href{mailto:gabel@mathematik.tu-darmstadt.de}{Fabian Gabel}
    \vspace*{1.0cm}

    15.10.2016 \\
    \vspace*{2cm}

    Betreuer: PD Dr. Robert Haller-Dintelmann

    \vspace*{.5cm}

    %Zweiter Gutachter: Name des zweiten Gutachters\\[2ex]
    \vspace*{\fill}
    \tiny{Version vom \today}
  \end{center}
\end{titlepage}
\vspace*{\fill}


% Inhaltsverzeichnis erstellen.
\tableofcontents

% Ab sofort werden Seitenzahlen in der Kopfzeile angezeigt.
\pagestyle{headings}

\chapter*{Einleitung}\index{Einleitung}
\addcontentsline{toc}{chapter}{Einleitung}\index{Einleitung}

\begin{addmargin}[2em]{2em}% 1em left, 2em right
  \textit{... as Sir Cyril Hinshelwood has observed ... fluid dynamicists
were divided into hydraulic engineers who observed things that
could not be explained and mathematicians who explained things
that could not be observed.} 
  \flushright(James Lighthill)
\end{addmargin}
\vspace{1.5cm}

Die akkurate Modellierung des Verhaltens \newton scher Fluide ist zentral für unzählige Anwendungen der Aerodynamik, Verbrennungsforschung oder chemischen Industrie.
Die Grundlage für die Modellierung bildet ein System partieller Differentialgleichungen, welches \navier\ und \stokes\ unabhängig voneinander in der ersten Hälfte des 19. Jahrhunderts einführten. 
Man bezeichnet dieses Gleichungssystem heute als \navier\hyp\stokes\hyp{}Gleichungen.
Sie umfassen im inkompressiblen dreidimensionalen Fall eine Gleichung zur Beschreibung der Massenerhaltung, die sogenannte Kontinuitätsgleichung, und für jede Raumrichtung eine Impulsgleichung:

\begin{align*}
  \div u &= 0, \\
  \frac{\partial u}{\partial t} - \nu \Delta u + u \cdot \nabla u + \nabla p &= f.
\end{align*}

Die Existenz und Eindeutigkeit klassischer glatter Lösungen dieses Gleichungssystems gehört nach dem \textsc{Clay}\hyp{}Institute zu einem der wichtigsten ungelösten mathematischen Probleme unseres Jahrtausends und ist daher in der Liste der Millennium Probleme zu finden.

In der Lösungstheorie der \navier\hyp\stokes\hyp{}Gleichungen aber auch partieller Differentialgleichungen im Allgemeinen hat es sich als nützlich erwiesen, auf der Suche nach klassischen Lösungen einen Umweg einzuschlagen und zunächst die Existenz sogenannter schwacher Lösungen nachzuweisen.
In einem zweiten Schritt wird dann überprüft, ob die gefundenen schwachen Lösungen über zusätzliche Regularitätseigenschaften verfügen, welche die Existenz einer Lösung im klassischen Sinne garantieren.

Als nützliche Werkzeuge für die Lösungstheorie haben sich dabei Zerlegungen geeigneter Funktionenräume erwiesen. Die Idee dazu fußt auf einer von \helmholtz\ erstmals 1870 formulierten Zerlegung eines Vektorfeldes in ein skalares Potential und ein Vektorpotential.
Im Rahmen der Suche nach schwachen Lösungen im \hilbert\hyp{}Raum $L^2(\Omega)^n$, wobei $\Omega$ ein Teilgebiet des $\R^n$, $n \geq 2$, bezeichne, findet auch das in dieser Arbeit vorgestellte Hilfsmittel Verwendung: die \helmholtz\hyp{}Zerlegung des $L^2$.
Dabei handelt es sich um eine orthogonale Zerlegung des Lösungsraumes, welche es unter anderem ermöglicht den neben dem gesuchten Geschwindigkeitsfeld $u$ unbekannten Druck $p$ vorerst aus dem Gleichungssystem zu eliminieren und so die Anzahl der Unbekannten zu reduzieren.
Die Anwendung der \helmholtz\hyp{}Zerlegung auf ein Element des Lösungsraumes zerlegt dieses in einen divergenzfreien und einen rotationsfreien Anteil.

Ziel dieser Arbeit ist es, aufbauend auf den im Grundstudium vermittelten Kenntnissen der Funktionentheorie, Integrationstheorie und Funktionalanalysis, die Existenz und Eindeutigkeit der \helmholtz\hyp{}Zerlegung auf dem Raum $L^2(\Omega)^n$ zu beweisen.

Das erste Kapitel stellt dazu die nötigen funktionalanalytischen Grundlagen bereit und führt in die in der Arbeit verwendete Notation ein.
Das Hauptaugenmerk dieses Kapitels liegt auf der Definition der schwachen Differenzierbarkeit und des Distributionsbegriffs.

Im zweiten Kapitel werden die zum Beweis der \helmholtz\hyp{}Zerlegung nötigen Hilfsaussagen bereitgestellt.
Das zentrale Resultat dieses Teils der Arbeit ist ein Kriterium, welches unter gewissen Bedingungen die schwache Lösbarkeit der Gradientengleichung $\nabla p = f$ sicherstellt.

Kapitel drei behandelt schließlich die \helmholtz\hyp{}Zerlegung $$L^2(\Omega)^n = L^2_\sigma(\Omega) \oplus G_2(\Omega)$$ und eine Charakterisierung dieser Zerlegung für den Fall, dass $\Omega = \R^n$ gilt.

Das letzte Kapitel fasst nochmals die zentralen Resultate zusammen und gibt einen Ausblick auf Problemstellungen, in denen die \helmholtz\hyp{}Zerlegung zum Einsatz kommt.

\chapter{Funktionalanalytische Grundlagen}
\label{cp:grundlagen}

Dieses Kapitel beschäftigt sich mit Definitionen und Eigenschaften der für die kommenden Kapitel zentralen Funktionenräume.
Die Notation orientiert sich an \cite{sohr2001navier}.
In Anlehnung an die für \hilbert\hyp{}Räume verbreitete Schreibweise $\langle \cdot, \cdot \rangle$ für das Skalarprodukt wollen wir nicht nur für den Raum $L^2$, sondern auch für alle anderen $L^q$\hyp{}Räume diese Symbolik verwenden.
Wir definieren also für geeignete Funktionen $f$ und $g$ den Ausdruck
$$
\langle f, g \rangle := \int f \cdot g \d x.
$$
Falls nicht anders vermerkt, bezeichnet $n$ immer eine von Null verschiedene natürliche Zahl.

\section{Glatte Funktionen und der Raum der Testfunktionen}
\label{subsec:smooth}

Ziel dieses Abschnittes ist es, die nötigen Begriffe und Definitionen im Zusammenhang mit glatten Funktionen bereitzustellen.
Im Folgenden bezeichne $\Omega \subseteq \R^n$ stets ein Gebiet, also eine nichtleere, offene und zusammenhängende Teilmenge.

Sei $k \in \N$. Mit $C^k(\Omega)$ bezeichnen wir die Menge aller Funktionen $u$, deren partielle Ableitungen $D^\alpha u$ für alle Multiindices $\alpha \in \N_0^n$ mit $0 \leq |\alpha| \leq k$ existieren und stetig sind. 
Mit
$$
  C^\infty(\Omega) := \bigcap_{k = 0}^\infty C^k(\Omega)
$$
bezeichnen wir den Raum der glatten Funktionen auf $\Omega$.

Im Kontext der Approximation von $L^q$-Funktionen spielt ein Funktionenraum eine wichtige Rolle: der Raum der glatten Funktionen mit kompaktem Träger
$$
C_0^\infty(\Omega) := \{ u \in C^\infty(\Omega) \mid \supp u \text{ kompakt und } \supp u \subseteq \Omega\}.
$$
Wir werden diesem Raum in Abschnitt \ref{subsec:distributionsSobolev} nochmals als Raum der Testfunktionen begegnen.
Zudem werden wir den Raum $C^\infty(\overline\Omega)$ aller Restriktionen $u|_{\overline\Omega}$ von Funktionen aus $C^\infty(\R^n)$ mit
\begin{displaymath}
  \sup_{\alpha  \in \N^n, x \in \R^n} |D^\alpha u(x)| < \infty 
\end{displaymath}
benötigen.
Aufgrund dieser Eigenschaft lässt sich der Raum $C^\infty(\overline\Omega)$ mit einer Norm ausstatten:
\begin{equation}
  \label{eq:cInftyNorm}
  \norm{u}_{C^\infty} := \norm{u}_{C^\infty(\overline\Omega)} := \sup_{|\alpha| \in \N^n, x \in \overline\Omega} |D^\alpha u(x)|.
\end{equation}

Alle eingeführten Räume lassen sich auf natürliche Weise zu Räumen von Vektor\-feldern verallgemeinern.
Man erhält so
\begin{align*}
  C^\infty(\Omega)^m &:= \{(u_1,\dots,u_m) \mid u_j \in C^\infty(\Omega), j = 1,\dots,m\}, \\
  C_0^\infty(\Omega)^m &:= \{(u_1,\dots,u_m) \mid u_j \in C_0^\infty(\Omega), j = 1,\dots,m\} \quad\text{und} \\
  C^\infty(\overline\Omega)^m &:= \{(u_1,\dots,u_m) \mid u_j \in C^\infty(\overline\Omega), j = 1,\dots,m\},
\end{align*}
wobei der letzte Vektorraum durch die Norm
$$
\norm{u}_{C^\infty} := \norm{u}_{C^\infty(\overline\Omega)}^m := \sup_{j = 1,\dots, m} \norm{u_j}_{C^\infty(\overline\Omega)}
$$
zu einem normierten Vektorraum wird.

In der Literatur findet man eine Reihe anderer Definitionen des Symbols $C^\infty(\overline\Omega)$.
Für beschränkte Gebiete stimmen sie jedoch alle überein, wie wir im Folgenden beweisen wollen.
Dazu benötigen wir ein Fortsetzungsresultat von \stein\ \cite[S.172, Proposition 2.2]{stein1970singular}.

\begin{prop}
  \label{prop:stein}
  Sei $f$ eine Funktion auf der abgeschlossenen Menge $F \subseteq \R^n$.
  Dann existiert eine Fortsetzung $\mathcal{E}_0(f)$ von $f$ auf $\R^n$.
  Ist $f$ zusätzlich stetig auf $F$, so ist die Fortsetzung $\mathcal{E}_0$ stetig auf $\R^n$ und glatt auf $F^c$.
\end{prop}

\newpage
\begin{lem}
  \label{lem:CInftyClosedOmega}
  Sei $\Omega \subseteq \R^n$ ein beschränktes Gebiet.
  Für eine Funktion $u \colon \overline\Omega \to \R$ schreiben wir:
  \begin{enumerate}[(a)]
    \item $u \in C^\infty(\overline\Omega)$,
      falls eine Funktion $f \in C^\infty(\R^n)$ mit $\sup_{\alpha \in \N^n, x \in \R^n} | D^\alpha f(x) | <  \infty$ existiert mit $u = f|_{\overline\Omega}$. \cite[S.23, I.3.1]{sohr2001navier}

    \item $u \in C^\infty_{\bullet}(\overline\Omega)$,
      falls eine offene Obermenge $\overline\Omega \subseteq U \subseteq \R^n$ und eine glatte Fortsetzung $\tilde u \colon U \to \R$ von $u$ existieren.

    \item $ u \in C^\infty_{\bullet\bullet}(\overline\Omega)$,
    falls $u|_\Omega \in C^\infty(\Omega)$ gilt und sich für alle Multiindices $\alpha \in \N^n$ die Funktion $D^\alpha u$ stetig auf $\overline\Omega$ fortsetzen lässt. \cite[S.10, 1.28]{adams2003sobolev}, \cite[S.35, II.1.3]{galdi2011navier}
  \end{enumerate}

  Dann gilt
  $$
    C^\infty(\overline\Omega) =
    C^\infty_{\bullet}(\overline\Omega) =
    C^\infty_{\bullet\bullet}(\overline\Omega).
  $$
\end{lem}

\begin{proof}
  %Zunächst ist $\overline\Omega$ als beschränkte abgeschlossene Teilmenge von $\R^n$ kompakt.
  Wir erhalten die Inklusion
  $$
  C^\infty(\overline\Omega) \subseteq C_{\bullet},
  $$ 
  da jedes $u\in C^\infty(\overline\Omega)$ bereits eine glatte Fortsetzung auf die offene Obermenge $\R^n$ besitzt.

  Die Inklusion $$C^\infty_{\bullet}(\overline\Omega) \subseteq C^\infty_{\bullet\bullet}(\overline\Omega)$$ folgt direkt aus der Definition.
  
  Die letzte Inklusion $$C^\infty_{\bullet\bullet}(\overline\Omega) \subseteq C^\infty(\overline\Omega)$$ erfordert mehr Arbeit.
  Sei dazu $u \in C^\infty_{\bullet\bullet}(\overline\Omega)$.
  %Es ist eine Eigenschaft von \hausdorff\hyp{}Räumen, dass stetige Abbildungen höchstens eine Fortsetzung auf den Abschluss ihres Definitionsberechs besitzen \cite[S.112, Lemma 4.2.4]{bartsch2015topologie}.
  %Wir wollen daher $u$ mit seiner Fortsetzung auf $\overline\Omega$ identifizieren.
  Das Fortsetzungsresultat von \stein\ \ref{prop:stein} erlaubt es, die stetige Funktion $u$ ausgehend von der abgeschlossenen Menge $\overline\Omega$ auf ganz $\R^n$ zu einer Funktion $\tilde u$ fortzusetzen.
  Diese Fortsetzung besitzt jedoch im Allgemeinen keine beschränkten Ableitungen wie in der Definition von $C^\infty(\overline\Omega)$ gefordert.
  Wir werden daher $\tilde u$ außerhalb von $\overline\Omega$ geeignet modifizieren.
  Da $\Omega$ als beschränkt vorausgesetzt wurde, existiert ein offener Ball $B_r(0)$ mit Radius $r > 0$ und 
  $$
   \Omega \subset\subset B_r(0).
  $$ 
  Wir multiplizieren nun die Fortsetzung $\tilde u$ mit einer glatten \emph{cut-off}\hyp{}Funktion $\psi$, für die
  $$ 
  \psi(x) = \begin{cases}
                          1 &\text{ für alle } x \in \overline\Omega \\
                          0 &\text{ für alle } x \in \R^n \setminus B_r(0)
             \end{cases}
  $$
  und folglich $\supp \psi \subseteq B_r(0)$ gilt.
  Eine solche Funktion erhält man zum Beispiel aus der Faltung der charakteristischen Funktion $\chi_{\overline\Omega}$ mit einem Glättungskern, wie der Beweis des \urysohn\hyp{}Lemmas für glatte Funktionen in \cite[S.88, Proposition 6.5]{folland2009ra} zeigt. 
  Für $f := \tilde u \cdot \psi$ erhalten wir
  $$
  \sup_{|\alpha| < \infty, x \in \R^n} | D^\alpha f |
   = \sup_{|\alpha| < \infty, x \in \overline B_r} | D^\alpha f |
   < \infty,
  $$
  da $\overline B_r$ kompakt ist.
  Aufgrund der Fortsetzungseigenschaft gilt weiterhin $f|_{\overline\Omega} = u$ und damit ist $f$ die für (a) gesuchte Fortsetzung.
\end{proof}

Zuletzt erwähnen wir den Untervektorraum der divergenzfreien glatten Vektorfelder
$$
C_{0,\sigma}^\infty(\Omega) := \{u \in C_0^\infty(\Omega)^n \mid \div u = 0\},
$$
 welcher in der schwachen Lösungstheorie der stationären inkompressiblen \navier\hyp{}\stokes\hyp{}Gleichungen als Testraum eingesetzt wird.

\section{Schwache Differenzierbarkeit}
\label{subsec:distributionsSobolev}

Wir wenden uns nun einer Verallgemeinerung des Differenzierbarkeitsbegriffs zu, welcher für die Lösungstheorie partieller Differentialgleichungen von zentraler Bedeutung ist: die \emph{schwache} Differenzierbarkeit.
Es bezeichne im Folgenden wieder $\Omega \subseteq \R^n$ ein Gebiet mit $n \geq 1$.

\subsection{Lokal integrierbare Funktionen}

Wir werden im Folgenden oft eine Obermenge der bezüglich des \lebesgue\hyp{}Maßes auf $\R^n$ integrierbaren Funktionen verwenden:
Für $1 \leq q \leq \infty$ schreiben wir
$$
  u \in L^q_{\loc}(\Omega)
$$
und nennen $u$ \emph{lokal integrierbar}, falls $u \in L^q(B)$ für alle offenen Bälle $B \subseteq \Omega$ mit $\overline B \subseteq \Omega$ gilt.

\begin{bem}
  Eine Funktion $u$ ist genau dann lokal integrierbar über $\Omega$, falls $u \in L^q(K)$ für alle Kompakta $K \subseteq \Omega$ gilt.

  Diese Aussage findet man ebenfalls in der Literatur zur Definition lokaler Integrierbarkeit.
  Tatsächlich ist sie äquivalent zur obigen Definition.
  Denn einerseits ist für alle Bälle $B$ auch $\overline B$ ein Kompaktum und die Aussage folgt aus der Inklusionsbeziehung $L^q(\overline B) \subseteq L^q(B)$.
  Andererseits lässt sich jedes Kompaktum $K$ durch endlich viele Bälle $B_1, \dots, B_n$ mit $\overline B_i \subseteq \Omega, i =1,\dots,m,$ überdecken. 
  Gilt nun $u \in L^q(B_i)$ für alle $i = 1, \dots, m$, so gilt insbesondere $u \in L^q(\bigcup_{i = 1}^m B_i)$ und die Umkehrung der Aussage folgt aus der Inklusionsbeziehung $L^q(\bigcup_{i = 1}^m B_i) \subseteq L^q(K)$.
\end{bem}

Des Weiteren schreiben wir
$$
u \in L^q_{\loc}(\overline\Omega),
$$
falls $u \in L^q(B \cap \Omega)$ für jeden offenen Ball $B \subseteq \R^n$ mit $B \cap \Omega \neq \emptyset$ gilt.
Zusammenfassend gilt also die folgende Inklusionsbeziehung:
$$
L^q(\Omega) \subseteq L^q_{\loc}(\overline\Omega) \subseteq L^q_{\loc}(\Omega).
$$

\subsection{Distributionen}

In Abschnitt \ref{subsec:smooth} hatten wir bereits den Raum der Testfunktionen $C_0^\infty(\Omega)$ kennengelernt.
Wir werden uns vor allem für seinen Dualraum, den Raum der stetigen Funktionale auf $C_0^\infty(\Omega)$, interessieren.
Um überhaupt über Stetigkeit von Funktionalen auf $C_0^\infty(\Omega)$ sprechen zu können, müssen wir jedoch zunächst eine Topologie auf dem besagten Raum festlegen.
Dabei beschränken wir uns auf den aus dieser Topologie resultierenden Stetigkeitsbegriff:
Ein lineares Funktional $F \colon C_0^\infty(\Omega) \to \R$ ist stetig, genau dann, wenn für jedes beschränkte Teilgebiet $G \subseteq \Omega$ mit $\overline G \subseteq \Omega$ ein $k \in \N_0$ und $C = C(F,G)$ existiert, sodass
$$
|F(\varphi)| \leq C \norm{\varphi}_{C^k(\overline G)}
$$
gilt, wobei $\norm{\cdot}_{C^k(\overline G)}$ gerade die in Gleichung (\ref{eq:cInftyNorm}) definierte Norm bezeichnet.
Die hinter diesem Stetigkeitsbegriff stehende Topologie macht $C_0^\infty(\Omega)$ zu einem lokalkonvexen Vektorraum.
Eine ausführlichere Beschreibung der beschriebenen Topologie findet sich in \cite[S.433f.]{werner2011fa}.
Den Raum $C_0^\infty(\Omega)'$ aller im obigen Sinne stetigen Funktionale
$$
F \colon C_0^\infty(\Omega) \to \R, \quad \varphi \mapsto F(\varphi) =: [F, \varphi],
$$
bezeichnen wir als den Raum der \emph{Distributionen}.
Hierbei stellt $[\cdot,\cdot]$ die duale Paarung auf $\Omega$ dar. 

Wir betrachten nun spezielle Distributionen.
Ist $f \in L^1_{\loc}$, so verstehen wir unter 
$$
f \mapsto \langle f, \varphi \rangle := \int_\Omega f \varphi \d x, \quad \varphi \in C_0^\infty(\Omega), 
$$
die zu $f$ gehörige \emph{reguläre} Distribution.
Dass diese Zuordnung sogar injektiv ist und somit zur Inklusionsbeziehung $L^1_{\loc}(\Omega) \subseteq C^\infty_0(\Omega)'$ führt, lässt sich in \cite[S.432, Beispiel (a)]{werner2011fa} nachlesen.

\subsection{Differentiation von Distributionen}

Distributionen ermöglichen es, den für die Analysis grundlegenden Begriff der Differenzierbarkeit zu verallgemeinern.
Ist $\alpha \in \N_0^n$ ein Multiindex und $D^\alpha \colon C_0^\infty(\Omega) \to C_0^\infty(\Omega)$ der zugehörige Ableitungsoperator, so definieren wir für eine Distribution $F \in C_0^\infty(\Omega)'$ ihre distributionelle Ableitung $D^\alpha F \in C_0^\infty(\Omega)'$ durch
$$
  [D^\alpha F, \varphi] := (-1)^{|\alpha|} [F, D^\alpha \varphi].
$$
Bis auf ein Vorzeichen stimmt also die distributionelle Ableitung mit der zum Ableitungsoperator dualen Abbildung $(D^\alpha)' \colon C_0^\infty(\Omega)' \to C_0^\infty(\Omega)'$ überein.
Dass die distributionelle Ableitung wohldefiniert und sogar schwach*-stetig ist, lässt sich in \cite[S.434, Lemma VIII.5.7]{werner2011fa} nachlesen.

Im Unterabschnitt zu glatten Funktionen hatten wir die dortigen Definitionen auf Produkträume ausgeweitet.
Wir versehen den so entstandenen Produktraum $C_0^\infty(\Omega)^m$ mit der Produkttopologie und wenden uns nun dem zugehörigen Distributionenraum $(C_0^\infty(\Omega)^m)'$ zu.

Wir betrachten dazu die Distribution
$$
  F = (F_1, \dots, F_m), \quad F_j \in C_0^\infty(\Omega)', \quad j = 1, \dots, m,
$$
und definieren für alle 
$$\varphi = (\varphi_1, \dots, \varphi_m) \in C_0^\infty(\Omega)^m$$
die Distribution
\begin{equation}
  \label{eq:pairingVector}
  [F, \varphi] := [F_1, \varphi_1] + \dots + [F_m, \varphi_m].
\end{equation}
Durch Betrachtung von Funktionen des Typs 
$$\psi^{(i)} = (0, \dots,0, \psi_i, 0,\dots,0) \in C_0^\infty(\Omega)^m, \quad i \in \{1,\dots,m\}, $$
erkennen wir, dass jedes bezüglich der Produkttopologie stetige Funktional auf $C_0^\infty(\Omega)^m$ diese Form besitzt.
Für den Raum der Distributionen auf den Testfunktionen $C_0^\infty(\Omega)^m$ gilt also
\begin{align*}
  (C_0^\infty(\Omega)')^m 
  &= (C_0^\infty(\Omega)^m)' \\
  &= \{(F_1,\dots,F_m) \mid F_j \in C_0^\infty(\Omega)', j = 1, \dots,m\}.
\end{align*}

Wir erinnern uns an den im Unterabschnitt zu glatten Funktionen definierten Funktionenraum $C_{0, \sigma}^\infty(\Omega) \subseteq C_0^\infty(\Omega)^n$ der divergenzfreien Funktionen.
Nach dem Fortsetzungssatz von \hahn\hyp{}\banach\ für lokalkonvexe Vektorräume (\cite[S.408, Satz VIII.2.8]{werner2011fa}) erhalten wir die stetigen Funktionale auf $C_{0,\sigma}^\infty$ gerade als Einschränkungen der stetigen Funktionale auf $C_0^\infty(\Omega)^n$. 
Es gilt somit
$$
C_{0,\sigma}^\infty(\Omega)' = \{F|_{C_{0,\sigma}^\infty(\Omega)} \mid F \in (C_0^\infty(\Omega)^n)'\}.
$$

Betrachten wir nun den \hilbert\hyp{}Raum $L^2(\Omega)^n$.
Wie Lemma \ref{lem:mollification} noch zeigen wird, liegen die glatten Funktionen mit kompaktem Träger dicht in $L^2(\Omega)$. 
Ebenso stellt der Produktraum $C_0^\infty(\Omega)^n$ einen dichten Unterraum von $L^2(\Omega)^n$ dar.
Um diese Tatsache zu imitieren, definieren wir den Unterraum 
$$
  L^2_\sigma(\Omega) 
  := \overline{C_{0,\sigma}^\infty(\Omega)}^{\norm{\cdot}_2}
  \subseteq L^2(\Omega)^n.
$$
Identifizieren wir nun jede Funktion $u \in L^2(\Omega)^n$ mit dem Funktional
$$
\langle u, \cdot \rangle \colon \varphi \mapsto \langle u, \varphi \rangle, \quad \varphi \in C_0^\infty(\Omega)^n,
$$
so erhalten wir die Einbettung 
$$
  L^2(\Omega)^n \subseteq (C_0^\infty(\Omega)^n)'.
$$
Ebenso lassen sich Funktionen $u \in L^2_\sigma(\Omega)$ mit der entsprechenden Einschränkung
$$
\langle u, \cdot \rangle \colon \varphi \mapsto \langle u, \varphi \rangle, \quad \varphi \in C_{0,\sigma}^\infty(\Omega)^n
$$
identifizieren, was die Einbettung
$$
  L^2_\sigma(\Omega) \subseteq (C_{0,\sigma}^\infty(\Omega)^n)'
$$
liefert.

\subsection{Sobolev-Räume}

Im Allgemeinen wird die Ableitung einer regulären Distribution nicht wieder regulär sein.
Wir definieren daher für alle $k \in \N$ und $1 \leq q \leq \infty$ den $L^q$\hyp{}\sobolev\hyp{}Raum $W^{k,q}(\Omega)$ der Ordnung $k$ durch
$$
W^{k,q}(\Omega) := \{ u \in L^q(\Omega) \mid D^\alpha u \in L^q(\Omega) \text{ für alle } |\alpha| \leq k\}.
$$
Für $u \in W^{k,q}(\Omega)$ bezeichnen wir $D^\alpha u \in L^q(\Omega)$ als \emph{schwache Ableitung} von $u$.
Hierbei identifizieren wir die reguläre Distribution $D^\alpha u$ immer direkt mit der korrespondierenden $L^q$\hyp{}Funktion.

Wir machen $W^{k,q}$ zu einem normierten Vektorraum durch die \sobolev\hyp{}Norm
$$
  \norm{u}_{W^{k,q}(\Omega)} 
  := \norm{u}_{W^{k,q}} 
  := \norm{u}_{k,q} 
  := \begin{cases} 
    \left(\sum_{|\alpha| \leq k} \norm{D^\alpha u}_q^q \right)^\frac{1}{q}\quad&\text{für } 1 \leq q < \infty \\
    \max_{|\alpha| \leq k} \norm{D^\alpha u}_\infty \quad&\text{für }q = \infty.
    \end{cases}
$$

\begin{lem}
  Für ein Gebiet $\Omega \subseteq \R^n$ mit $n \geq 1$ und $1 \leq q \leq \infty$ ist $W^{k,q}(\Omega)$ ausgestattet mit der \sobolev\hyp{}Norm ein \banach\hyp{}Raum.
\end{lem}

\begin{proof}
  Sei $(u_j)_{j \in \N}$ eine \cauchy\hyp{}Folge in $W^{k,q}(\Omega)$.
  Nach Definition ist dann auch $(D^\alpha u_j)_{j \in \N}$ für alle $|\alpha| \leq k$ eine \cauchy\hyp{}Folge in $L^q$.
  Folglich existieren für alle $|\alpha| \leq k$ Funktionen $u_\alpha \in L^q$ mit 
  $$
  \lim_{j \to \infty} \norm{D^\alpha u_j - u_\alpha}_q = 0.
  $$
  Sei $u_0 := u_{(0,\dots,0)}$.
  Wir behaupten nun, dass im distributionellen Sinne die Identität $u_\alpha = D^\alpha u_0$ für alle $\alpha \in \N^n$ gilt.
  Wir zeigen dazu, dass für alle $\varphi \in C_0^\infty(\Omega)$ die Identität
  $$
  [u_\alpha, \varphi] = [D^\alpha u_0, \varphi]
  $$
  gilt.
  Dazu halten wir zunächst fest, dass mit der \hoelder\hyp{}Ungleichung
  $$
  \int_\Omega |(D^\alpha u_j - u_\alpha)| \varphi \d x
  \leq \norm{(D^\alpha u_j - u_\alpha)}_q \norm{\varphi}_{q'} \to 0, \quad j \to \infty
  $$
  sowie
  $$
  \int_\Omega |(u_j - u_0) D^\alpha \varphi| \d x
  \leq \norm{(u_j - u_0)}_q \norm{D^\alpha \varphi}_{q'} \to 0, \quad j \to \infty
  $$
  folgen, wobei $q'$ den zu $q$ konjugierten Exponenten bezeichne.
  Hieraus folgen die Identitäten
  $$
  \lim_{j \to \infty} [D^\alpha u_j, \varphi] = \lim_{j \to \infty} \int_\Omega D^\alpha u_j \varphi \d x =  \int_\Omega u_\alpha \varphi \d x =  [u_\alpha, \varphi] 
  $$
  und
  $$
    \lim_{j \to \infty} [u_j, D^\alpha\varphi] = \lim_{j \to \infty} \int_\Omega u_j D^\alpha\varphi \d x =  \int_\Omega u_0 D^\alpha \varphi \d x = [u_0, D^\alpha \varphi].
  $$
  Dies impliziert schließlich
  $$ 
  [u_\alpha, \varphi] 
  = \lim_{j \to \infty} [D^\alpha u_j, \varphi] \\
  = \lim_{j \to \infty} (-1)^{|\alpha|}[u_j, D^\alpha \varphi] \\
  =  (-1)^{|\alpha|}[u_0, D^\alpha \varphi]
  = [D^\alpha u_0, \varphi].
  $$
  Daraus folgt jedoch gerade die Behauptung $D^\alpha u_0 = u_\alpha$.
  Insgesamt ergibt sich also $\lim_{j \to \infty} \norm{u_j - u_0}_{k,q} =  0$, was zu beweisen war.
\end{proof}

Wir definieren den Unterraum
$$
  W_0^{k,q}(\Omega) := \overline{C_0^\infty(\Omega)}^{\norm{\cdot}_{k,q}},
$$
welcher die glatten Funktionen mit kompaktem Träger als dichte Teilmenge bezüglich der \sobolev\hyp{}Norm enthält.

Bisher hatten wir nur \sobolev\hyp{}Räume positiver Ordnung betrachtet.
Wir definieren für $1 < q < \infty$ die \sobolev\hyp{}Räume negativer Ordnung 
$$
  W^{-k,q}(\Omega) := W_0^{k,q'}(\Omega)'.
$$
Wie für den Dualraum eines normierten Raumes üblich, ist $W^{-k,q}(\Omega)$ mit Operatornorm 
$$
\norm{F}_{W^{-k,q}(\Omega)} := \norm{F}_{W^{-k,q}} := \norm{F}_{-k,q} := \sup_{0 \neq \varphi \in C_0^\infty(\Omega)} \frac{|[F,\varphi]|}{\norm{\varphi}_{k,q'}}
$$
ein \banach\hyp{}Raum.
Wir beachten, dass die Operatornorm bereits durch die bezüglich \sobolev\hyp{}Norm in $W^{k,q}(\Omega)$ dichte Teilmenge $C_0^\infty(\Omega)$ eindeutig festgelegt ist.

Es ist möglich einen \sobolev\hyp{}Raum $W^{k,q}(\Omega)$ für $1 < q < \infty$ als abgeschlossenen Teilraum eines $L^q$\hyp{}Raumes zu realisieren \cite[S.61, 3.5]{adams2003sobolev}.
Als solcher ist er insbesondere reflexiv.
Wir halten diese wichtige Eigenschaft in einem Lemma fest.

\begin{lem}
  Für ein Gebiet $\Omega \subseteq \R^n$ mit $n \geq 1$ und $1 < q < \infty$ ist $W^{k,q} = W^{k,q}(\Omega)$ ausgestattet mit der \sobolev\hyp{}Norm ein reflexiver Raum.
\end{lem}

Aus der Reflexivität der \sobolev\hyp{}Räume leiten wir für $1 < q < \infty$ die folgende Identität ab:
$$
  W_0^{k,q'}(\Omega) = W^{-k,q}(\Omega)'.
$$
Wir können also jede Funktion $u \in W_0^{k,q'}(\Omega)$ mit dem Funktional
$$
[\cdot, u] \colon F \mapsto [F, u], \quad F \in W^{-1,q}(\Omega)
$$
identifizieren.

%Wir definieren nun die $W^{k,q}_{\loc}$\hyp{}Räume.
Wir bezeichnen mit $W^{k,q}_{\loc}(\Omega)$ den Raum aller Funktionen $u$ mit $D^\alpha u \in L^q_{\loc}(\Omega)$ für alle Multiindices $\alpha \in \N_0^n$ mit $|\alpha| \leq k$.
Des Weiteren definieren wir $W^{k,q}(\overline\Omega)$ als den Raum aller Funktionen $u$ mit $D^\alpha u \in L^q_{\loc}(\overline\Omega)$.
Nun zu den \sobolev\hyp{}Räumen negativer Ordnung. 
Wir definieren den Raum $W^{-k,q}_{\loc}(\Omega)$ als den Raum aller Distributionen
$$
F \colon \varphi \mapsto [F,\varphi], \quad \varphi \in C_0^\infty(\Omega),
$$
sodass für die Operatornorm die Ungleichung
\begin{equation}
  \label{eq:operatornorm}
\norm{F}_{-k,q} := \sup_{0 \neq \varphi \in C_0^\infty(\Omega_0)}  \frac{|[F,\varphi]|}{\norm{\varphi}_{k,q'}} < \infty
\end{equation}
für alle beschränkten Teilgebiete $\Omega_0 \subseteq \Omega$ mit $\overline\Omega_0 \subseteq \Omega$ gilt.

Wie schon im Unterabschnitt zu glatten Funktionen vorgestellt, können wir auch die Definition der \sobolev\hyp{}Räume auf Produkträume ausdehnen. 
Wir definieren hierzu
$$
W^{k,q}(\Omega)^m := \{ (u_1,\dots,u_m) \mid u_j \in W^{k,q}(\Omega), \quad j = 1,\dots, m\}
$$
und versehen den so entstandenen Vektorraum mit der Norm
$$
\norm{u}_{W^{k,q}(\Omega)^m} 
:= \norm{u}_{k,q}
:= \left( \sum_{j = 1}^m \norm{u_j}_{k,q}^q \right)^{\frac{1}{q}},
$$
wodurch dieser zu einem \banach\hyp{}Raum wird.


%\newpage
\section{Eigenschaften der Faltung}
\label{subsec:mollification}

Aus der Integrationstheorie ist bekannt, dass sich $L^q$-Funktionen durch Faltungen mit Glättungskernen $\mathcal{F}_\varepsilon$ approximieren lassen.
Für unsere Zwecke wird es dafür ausreichend sein den Standardkern
$$
\mathcal{F}(x) := \begin{cases}
                     c \cdot \exp\left(-\frac{1}{1 - \norm{x}}\right), \quad &\norm{x} < 1 \\
                     0, \quad  &\text{sonst}
                  \end{cases} \quad,\quad x \in \R^n
$$
zu betrachten und
$$
\mathcal{F}_\varepsilon(x) := \frac{1}{\varepsilon^n} \mathcal{F}\left(\frac{x}{\varepsilon}\right)
$$ 
zu setzen.
Hierbei sei $c$ so gewählt, dass $\int_{B_1(0)} \mathcal{F} \d x = 1$ gilt.

Folgendes Lemma fasst nochmals die Approximationseigenschaft von Glättungskernen zusammen. Ein Beweis findet sich in \cite[S.36, Theorem 2.29(c)]{adams2003sobolev}.
\begin{lem}
  \label{lem:mollification}
  Sei $\Omega \subseteq \R^n$ mit $n \geq 1$ ein Gebiet und $1 \leq q < \infty$ sowie $\varepsilon > 0$.
  Für alle $u \in L^q(\Omega)$ gilt dann
  $$
  \norm{(\mathcal{F}_\varepsilon \ast u)}_{L^q(\Omega)} \leq  \norm{u}_{L^q(\Omega)}
  $$
  und
  \begin{displaymath}
    \lim_{\varepsilon \to 0} \norm{(\mathcal{F}_\varepsilon \ast u) - u}_{L^q(\Omega)} = 0. 
  \end{displaymath}
\end{lem}

Sei nun $\Omega_0 \subseteq \overline\Omega_0 \subseteq \Omega$ ein beschränktes Teilgebiet und
\begin{equation}
  \label{eq:convolutionEpsilon}
  0 < \varepsilon < \dist(\Omega_0, \partial\Omega).
\end{equation}
Zusätzlich sei $u \in L^1_{\loc}(\Omega)$.
Setzt man $u(x) := 0$ für alle $x \in \R^n \setminus \Omega$, so folgt $u \in L^1_{\loc}(\R^n)$.
Wie \cite[S.171, Theorem 6.30(b)]{rudin1991fa} zeigt, lassen sich aus der Integrationstheorie bekannte Eigenschaften der Glättungen auch auf Faltungen 
$$
u^\varepsilon := \mathcal{F}_\varepsilon \ast u
$$
von Glättungskernen $\mathcal{F}_\varepsilon \ast u$ mit Distributionen $u \in C_0^\infty(\Omega)'$ verallgemeinern.
Einerseits zählt dazu, dass $u^\varepsilon \in C^\infty(\R^n)$ gilt.
Da $\Omega_0$ als beschränkt vorausgesetzt wurde, gilt nach Lemma \ref{lem:CInftyClosedOmega}(b$\Rightarrow$a) sogar $u^\varepsilon \in C^\infty(\overline\Omega_0)$.
Andererseits ist für alle Multiindices $\alpha \in \N_0^n$ und alle $x \in \Omega_0$ die folgende Gleichung gültig:
\begin{equation}
  \label{eq:convolutionDiff}
  (D^\alpha u^\varepsilon)(x) 
  = (\mathcal{F}_\varepsilon \ast (D^\alpha u))(x)
  = ( (D^\alpha \mathcal{F}_\varepsilon) \ast u)(x), \quad x \in \Omega_0.
\end{equation}

\begin{lem}
\label{lem:nablauconstant}
  Sei $\Omega \subseteq \R^n$ mit $n \geq 1$ ein Gebiet und $1 \leq q < \infty$.
  Ist $u \in L^1_{\loc}(\Omega)$ und gilt $\nabla u = 0$ im distributionellen Sinne, dann ist $u$ konstant.
\end{lem}

\begin{proof}
  Für alle beschränkten Teilgebiete $\Omega_0 \subseteq \Omega$ und $\varepsilon$ wie in Ungleichung (\ref{eq:convolutionEpsilon}) folgern wir unter Verwendung der Gleichung (\ref{eq:convolutionDiff})
  \begin{align*}
    \nabla u^\varepsilon
    &= (D_1 (\mathcal{F}_\varepsilon \ast u), \dots, D_n (\mathcal{F}_\varepsilon \ast u))^T \\
    &= ( \mathcal{F}_\varepsilon \ast D_1 u, \dots, \mathcal{F}_\varepsilon \ast D_n u)^T \\
    &= (0, \dots, 0)^T.
  \end{align*}
  Es ist $u^\varepsilon$ eine glatte Funktion.
  Zudem ist $\Omega_0$ als offene zusammenhängende Teilmenge von $\R^n$ insbesondere wegzusammenhängend.
  Ausgehend von der Integraldarstellung des Funktionszuwachses \cite[S.57, (11)]{koenigsberger2004ana2} gilt $u_\varepsilon = C_\varepsilon$ auf ganz $\Omega_0$.
  Mit Lemma \ref{lem:mollification} folgt nun die Netzkonvergenz $C_\varepsilon \to C$ für $\varepsilon \to 0$ auf $\Omega_0$.
  Wir greifen auf eine im Abschnitt \ref{sec:lipschitzDomains} zu \lipschitz\hyp{}Gebieten eingeführte Konstruktion vor:
  Aufgrund von Lemma \ref{lem:lipschitzExhaust} lässt sich das Gebiet $\Omega$ von einer Folge $(\Omega_j)_{j \in \N}$ offener beschränkter \lipschitz\hyp{}Gebiete ausschöpfen. 
  Durch Anwendung des für das Gebiet $\Omega_0$ beschriebenen Arguments auf die Folgenglieder $\Omega_j$ erhalten wir somit $u = C$ auf ganz $\Omega$. 
  Hieraus folgt die Behauptung.
\end{proof}





%\nocite{*}
\bibliographystyle{alphadin}
\bibliography{bachelorthesis_fabian_gabel}
%\addcontentsline{toc}{chapter}{Literaturverzeichnis}
\addcontentsline{toc}{chapter}{Literaturverzeichnis}


\renewcommand{\glossarysection}[2][\theglstoctitle]{%
    \def\theglstoctitle{#2}%
      \setcounter{secnumdepth}{-1}%
        \chapter[#1]{#2}%
}
\printglossary[title=Symbolverzeichnis,toctitle=Symbolverzeichnis]
      
\end{document}
